% PAPER STATUS (2026-07-24): authoritative efficiency ledger for StreamWeave.
% Public headline: resource-matched end-to-end prompt-group throughput
% 2.78 -> 4.58 groups/s (1.64x), or 46.0 -> 28.0 s per 128 groups.
% Main-text mechanism evidence: active-GPU count distribution from 15-second
% training-interval telemetry (SM-active threshold 20%; sensitivity 10--50%).
% Exact hardware and topology are Appendix-only details.
% Do not restore historical CISPO-arm claims such as 1.54x, 54% -> 4.6%
% trainer idle, 70% tail wait, 82% rollouter utilization, or MFU headlines.

Appendix B. 실행 효율: complete-group learning을 다시 직렬화하지 않기

B.0 측정 대상과 공통 작업 단위

이 부록의 질문은 단순하다. StreamWeave가 group-dependent learning을 단순화하지 않은 채,
complete-group decision이 요구하는 기다림을 전체 학습 pipeline의 직렬화로 전파하지 않는가?
이를 측정하기 위해 synchronous implementation과 StreamWeave를 동일한 8$\times$B200
hardware budget에서 비교한다. 전자는 colocated synchronous execution을, 후자는 trainer와
rollout generator를 분리한 fully-asynchronous execution을 사용한다. 정확한 run은 각각
\texttt{v96fvd0p}와 \texttt{oki4kv8u}다.

공통 작업 단위는 learner row나 cycle이 아니라 \emph{prompt group}이다. Policy source로
routing된 group은 8개의 learner row를 만들지만 expert source로 routing된 group은 하나의
row를 만들기 때문에, rows/s나 steps/s는 source mixture에 따라 달라진다. 반면 prompt group은
source가 결정되기 전부터 두 시스템에 공통이며 모든 group은 8개의 rollout attempt를 생성한다.
로그의 \texttt{hpt/onpolicy\_num\_groups}는 \texttt{group\_uid}를 중복 제거해 실제 소비된
prompt group을 센다. 따라서 학습 판단의 의미 단위와 실행 효율의 측정 단위가 일치한다.

Cycle 수는 직접 비교하지 않는다. Synchronous cycle은 128 groups를 소비하지만 StreamWeave
cycle은 261--489 groups, 평균 453.5 groups를 소비한다. 처리량은 cycle별 비율의 산술평균이
아니라 다음 work-weighted aggregate로 계산한다.

\[
\operatorname{Throughput}
=
\frac{\sum_i G_i}{\sum_i T_i},
\]

여기서 \(G_i\)는 cycle \(i\)에서 소비한 고유 prompt group 수이고, \(T_i\)는 evaluation을
제외한 training-loop time이다. 두 run의 W\&B full history를 직접 합산했으며 UI smoothing이나
downsampling된 값을 사용하지 않았다.

\begin{table}[h]
\centering
\begin{tabular}{lrrrrr}
\hline
System & Cycles & Groups & Time (s) & Groups/s & 128-group time (s) \\
\hline
Synchronous & 104 & 13{,}312 & 4{,}780.2 & 2.78 & 46.0 \\
StreamWeave & 190 & 86{,}174 & 18{,}828.4 & \textbf{4.58} & \textbf{28.0} \\
\hline
\end{tabular}
\caption{동일한 hardware budget에서 측정한 non-evaluation training-loop throughput.
StreamWeave는 prompt-group throughput을 1.64$\times$ 높인다. 128-group time은 같은
aggregate throughput의 역수를 환산한 값이며 별개의 증거가 아니다.}
\label{tab:streamweave-efficiency}
\end{table}

이 비교는 동일 자원에서 관찰한 end-to-end prompt-group throughput이다. Prompt group 수가
동일하더라도 생성 token 수가 항상 같다는 뜻은 아니므로, 이를 architecture-isolated speedup이나
동일 token workload 비교로 부르지 않는다. 또한 품질 threshold에 도달하는 시간을 맞춘 분석이
아니므로 time-to-quality 주장으로 확대하지 않는다.

B.1 Critical path가 어떻게 달라지는가

Synchronous implementation에서는 complete-group rollout generation이 serialized training
loop 안에 있다. Full history에서 generation은 non-evaluation training-loop time의 54.7\%를
차지한다. 128 prompt groups를 기준으로 generation에만 25.1초가 들고, 뒤이어 실행되는 학습까지
포함하면 전체가 46.0초다.

StreamWeave는 complete group 자체를 없애지 않는다. Trajectory attempt를 독립적으로 실행하고,
source decision에 필요한 시점에만 group을 복원함으로써 generation을 learner의 serialized
critical path 밖에서 계속 수행한다. 그 결과 generation과 training을 모두 포함한 전체 pipeline이
동일한 128-group work를 28.0초에 처리한다. 이는 synchronous implementation이 generation에만
쓰는 25.1초보다 2.8초 길 뿐이다. StreamWeave가 학습 계산을 줄인 것이 아니라, 동기식에서
generation 뒤에 직렬로 놓이던 계산의 대부분을 generation과 겹쳤다는 것이 첫 번째 효과다.

그러나 overlap만으로는 관찰된 critical path의 산술이 모두 닫히지 않는다. Synchronous
implementation은 generation에 8개 GPU를 사용하므로 유효 generation rate는

\[
r_{\mathrm{sync}}
=
\frac{128}{25.13\times 8}
=
0.637
\quad\text{groups/(GPU$\cdot$s)}.
\]

StreamWeave의 6-GPU rollouter가 전체 pipeline의 28.0초 처리율을 유지하려면, generation과
training을 모두 포함한 시간을 보수적인 분모로 사용해도 공급률이

\[
r_{\mathrm{StreamWeave}}
\ge
\frac{128}{27.97\times 6}
=
0.763
\quad\text{groups/(GPU$\cdot$s)}
\]

여야 한다. 이는 synchronous generation보다 약 1.20$\times$ 높은 유효 rate다. 동일한 GPU당
rate였다면 6-GPU generation만 약 33.5초가 필요했을 것이다. Queue cap 384 groups는 전체
86,174-group history의 0.45\%에 불과하므로 초기 backlog 소모로도 이 잔차를 설명할 수 없다.
따라서 두 번째 효과는 attempt-level scheduling이 complete-group tail을 기다리며 잃던 rollout-side
capacity를 회수한 것이다. 15초 간격의 training telemetry에서도 같은 실행 패턴이 직접 관찰된다.
SM activity가 20\%를 넘는 GPU가 하나도 없는 interval은 synchronous execution에서 27.9\%였지만
StreamWeave에서는 4.7\%였고, interval당 해당 GPU 수의 평균은 5.40에서 6.92로 증가했다. 이 방향은
threshold를 10--50\%로 바꾸어도 유지된다. 이 분포는 complete-group waiting이 전역적인 낮은
activity로 노출되는 빈도를 줄이고 더 많은 GPU에서 동시 작업을 유지한다는 기전 증거다. 양쪽 첫
training cycle을 제외해도 각각 24.9\% 대 4.7\%, 평균 5.63 대 6.93으로 같은 방향이 남는다.

이 계산은 speedup을 두 개의 독립적인 수치로 분해하는 것이 아니라, 하나의 1.64$\times$ end-to-end
결과가 overlap만으로 설명되지 않음을 보이는 consistency check다. Prompt group당 generated token
수, batching, colocated contention을 맞춘 비교가 아니므로 약 20\%를 barrier 제거 하나의
architecture-isolated causal effect로 부르지 않는다. Partial-rollout recovery도 불필요한 prefix
재생성을 줄일 수 있지만, 현행 로그는 그 기여율을 별도로 분리하지 않는다.

StreamWeave learner에서 ready group을 인출하고 역직렬화하며 조립하는 시간은 startup을 제외한
learner-side loop의 3.25\%다. 그러나 이 지표는 synchronous generation share와 동일한 low-level
operation을 재는 타이머가 아니다. 따라서 이를 \(54.7\%\to3.25\%\)의 stall 또는 trainer-idle
감소로 쓰지 않는다. 허용되는 구조적 해석은 다음과 같다.

\begin{quote}
Synchronous execution에서는 generation이 trainer의 critical path 안에 있지만,
StreamWeave에서는 generation이 그 경로 밖에서 진행되고 learner에는 완성된 group을 넘기는
짧은 interface만 남는다.
\end{quote}

B.2 서로 다른 run 길이에 대한 민감도

두 run의 cycle 수와 총 관찰 길이가 다르다는 사실은 Table~\ref{tab:streamweave-efficiency}의
aggregate estimator를 왜곡하지 않는다. 이를 직접 확인하기 위해 synchronous run의 전체 작업량과
동일한 13{,}312 groups만 StreamWeave history에서 잘라 다시 계산했다.

\begin{table}[h]
\centering
\begin{tabular}{lrrr}
\hline
Window & Time for 13{,}312 groups (s) & Groups/s & Relative throughput \\
\hline
Synchronous full history & 4{,}780.2 & 2.78 & 1.00$\times$ \\
StreamWeave first equal-work window & 2{,}856.3 & 4.66 & \textbf{1.67$\times$} \\
StreamWeave last equal-work window & 2{,}716.0 & 4.90 & \textbf{1.76$\times$} \\
\hline
\end{tabular}
\caption{동일 prompt-group budget으로 제한한 민감도 분석. 더 긴 StreamWeave history가
headline 1.64$\times$를 부풀리지 않았음을 보인다.}
\label{tab:equal-work-sensitivity}
\end{table}

StreamWeave history를 13{,}312-group 크기의 연속 구간으로 나누면 각 구간의 처리량은
4.66, 4.67, 3.48, 4.85, 5.16, 4.98 groups/s다. 가장 느린 구간도 synchronous throughput보다
약 1.25$\times$ 높다. 첫 cycle을 양쪽에서 제외하면 1.59$\times$, 알려진 parameter-sync
transient를 제외하면 1.65$\times$이며, 마지막 quarter도 4.95 groups/s를 유지한다. 따라서
1.64$\times$는 특정 peak나 짧은 구간을 선택한 결과가 아니다.

Cycle들은 시간적으로 연속된 시스템 관측이며 독립 표본이 아니다. 따라서 cycle을 i.i.d. sample로
간주한 표준오차나 유의확률은 보고하지 않는다. 본문의 headline은 full-history aggregate 하나로
고정하고, 동일 작업량 분석은 unequal run length에 대한 Appendix 방어로만 사용한다.

B.3 Boundedness, utilization, and cost ledger

Headline throughput에는 StreamWeave가 실제로 지불한 시스템 비용이 포함되어 있다. 190 cycles의
parameter synchronization은 median 2.15초, p95 2.79초였고, 약 46초, 104초, 200초의 세 transient도
full-history 1.64$\times$에 포함된다. Queue cap은 384 groups이며 startup 이후 collection
boundary에서의 occupancy는 약 58--75\%, 중앙값 70\%였다. 이는 해당 관측 경계에서 backlog가
유계였다는 근거이며, 모든 순간의 queue occupancy나 무조건적인 non-starvation을 뜻하지 않는다.

Carryover된 group의 기록상 discard는 전체 run에서 0이었다. 이는 carryover path에 대한 관측 결과일
뿐, 모든 종류의 producer-side drop이나 보편적 zero-waste를 주장하지 않는다. Partial rollout,
queue configuration, parameter-refresh protocol, framework-specific row alignment는 end-to-end
realization을 설명하는 보조 자산이며 StreamWeave의 독립 novelty로 올리지 않는다.

Source mixture도 속도 향상을 singleton expert row로 설명하지 못한다. StreamWeave의 RL-group
비율은 77.3\%로 synchronous run의 52.4\%보다 높았다. 다만 두 run의 protocol 차이가 있으므로
이를 정규화된 fairness proof로 사용하지 않고 방향성 있는 방어 근거로만 둔다. Training interval에
맞춘 GPU telemetry에서 전체 GPU busy는 StreamWeave 84.6\%, synchronous 65.5\%였고, trainer와
rollout partition이 동시에 50\% 이상 active였던 표본은 76.3\%였다. 이 aggregate 값과 per-GPU
heatmap은 Appendix 보조 증거로만 사용한다. 본문 Figure~3은 더 직접적으로 읽히는 active-GPU count
distribution만 사용하며, 15초 interval을 독립 실험 반복으로 취급하지 않는다. 또한 `zero active'는
어떤 GPU도 선언한 SM threshold를 넘지 않았다는 뜻이며 idle, stall, 0\% utilization과 동일하지 않다.

B.4 공개 주장 규율

본문의 실행 효율 주장은 다음 세 단계로 제한한다.

\begin{enumerate}
  \item \textbf{Metric.} Routing 이전의 공통 단위인 prompt group으로 work를 정의한다.
  \item \textbf{Mechanism.} Complete-group generation을 serialized critical path에서 분리하되,
        complete-group decision 자체는 유지한다. 이 분리는 generation과 training을 겹칠 뿐 아니라,
        group-tail waiting에 묶이던 rollout-side capacity도 회수한다. Active-GPU count distribution은
        전역적으로 낮은 activity의 감소와 더 지속적인 concurrent work를 직접 보여준다.
  \item \textbf{Payoff.} Full-history throughput은 2.78에서 4.58 groups/s로 증가하고,
        128-group-equivalent time은 46.0초에서 28.0초로 줄어든다. 이는 하나의 1.64$\times$
        end-to-end result다.
\end{enumerate}

가장 강한 해석은 더 빠른 async framework가 아니라 다음 architecture claim이다.

\begin{quote}
StreamWeave는 group-dependent learning을 단순화하지 않는다. Group completion이 필요한 지점만
기다리고, 그 기다림이 pipeline 전체의 직렬화나 rollout-side capacity loss로 전파되는 것을 막는다.
\end{quote}

과거 CISPO-arm에서 나온 1.54$\times$, \(54\%\to4.6\%\) trainer idle, 70\% tail wait,
82\% rollouter utilization, MFU headline, response-length-conditioned 41\% 증가는 현재 공개
근거로 사용하지 않는다. 특히 synchronous generation share와 asynchronous learner-side
acquisition/assembly share를 하나의 stall 감소율로 합치지 않는다.

% Measurement sources:
%   Synchronous: W&B run v96fvd0p, 104 cycles.
%   StreamWeave main: W&B run oki4kv8u, 190 cycles.
%   Numerator: sum of deduplicated hpt/onpolicy_num_groups.
%   Denominator: sum of timing_s/step over non-evaluation training intervals.
%   Code definition: verl/trainer/ppo/metric_utils.py, unique group_uid accounting.
%   Frozen snapshot and refresh path:
%     figures/execution_efficiency/data/verified_snapshot.json
%     figures/execution_efficiency/scripts/refresh_execution_efficiency.py
